\documentclass[11pt,a4paper]{article}
\usepackage[margin=1in]{geometry}
\usepackage{times}
\usepackage{latexsym}
\usepackage{graphicx}
\usepackage{amsmath}
\usepackage{amssymb}
\usepackage{booktabs}
\usepackage{multirow}
\usepackage{xcolor}
\usepackage{url}
\usepackage{tabularx}
\usepackage{microtype}
\usepackage[colorlinks=true,citecolor=blue,urlcolor=blue]{hyperref}

\title{Beyond Aggregate Benchmarks: Per-Example Analysis Exposes\\Architecture-Specific Behavior in Hybrid Language Models}

\author{Anonymous}

\date{}

\begin{document}
\maketitle

\begin{abstract}
Hybrid language models that combine linear recurrence with full attention are evaluated almost exclusively via aggregate benchmark scores.
We argue that this practice obscures architecture-specific behavior at the individual-example level.
We propose two per-example divergence metrics---pairwise Jaccard overlap of affected sample sets and net-effect asymmetry---and apply them to 30 hybrid models spanning three recurrence mechanisms (GatedDeltaNet, RetNet, HGRN2), five attention ratios, and two scales from the Wang et al.\ model zoo.
Our analysis reveals that even when aggregate cross-architecture variance decreases modestly with added attention (e.g., 38\% standard-deviation reduction on LAMBADA, $\rho=-0.60$, $p=0.117$), the underlying per-example effects are largely non-overlapping across architectures (mean pairwise Jaccard = 0.09, 95\% CI [0.03, 0.19]).
Crucially, no benchmark reaches conventional significance after multiple-comparison correction, and all aggregate convergence claims depend on a single confounded architecture (GatedDeltaNet).
These findings demonstrate that aggregate metrics can create misleading impressions of architectural equivalence, and motivate per-example analysis as a necessary complement to leaderboard evaluation.
We release our analysis framework and provide practitioner recommendations for architecture selection beyond aggregate scores.
\end{abstract}

\section{Introduction}
\label{sec:intro}

Hybrid language models that interleave linear recurrent layers with full attention layers have emerged as a promising design paradigm, offering sub-quadratic inference cost while retaining much of the expressivity of Transformers \cite{dao2024transformers,gu2024mamba}.
Recent work by Wang et al.\ provides the most comprehensive controlled comparison to date: a model zoo of 72 checkpoints spanning three linear recurrent architectures---GatedDeltaNet (GDN), RetNet, and HGRN2---at two parameter scales (340M and 1.3B) with systematically varied attention ratios.

When evaluated on standard benchmarks, these models exhibit a striking pattern: as attention layers are added, aggregate performance gaps across architectures narrow.
Wang et al.\ qualitatively note this convergence trend but do not analyze its per-example structure.
This leaves open a fundamental question: does aggregate convergence imply that architectures behave similarly on individual examples, or does it mask architecture-specific effects?

We address this question by introducing a per-example analysis framework with two core metrics:
\begin{enumerate}
    \item \textbf{Pairwise Jaccard overlap}: measuring whether attention helps or hurts the same examples across architectures.
    \item \textbf{Net-effect asymmetry}: quantifying whether attention's overall direction (beneficial vs.\ detrimental) differs across architectures on a per-example basis.
\end{enumerate}

Applying this framework to six benchmarks reveals that aggregate convergence metrics are fragile, non-significant after correction for multiple comparisons, and dependent on a confounded data point (GDN).
More importantly, even the modest aggregate convergence that does exist masks substantial per-example heterogeneity: architectures converge to similar aggregate scores by affecting largely disjoint sets of examples.
This finding has direct implications for practitioners: architecture selection based solely on leaderboard scores ignores systematic differences in which inputs each architecture handles well.

Our contributions are:
\begin{itemize}
    \item A per-example divergence framework (Jaccard overlap, net-effect asymmetry) for comparing hybrid architectures beyond aggregate metrics.
    \item A systematic case study on 30 models from the Wang et al.\ zoo demonstrating that aggregate metrics obscure per-example heterogeneity.
    \item Robustness analysis showing that aggregate convergence claims fail conventional significance thresholds and depend on a single confounded architecture.
    \item Practitioner recommendations for architecture selection informed by per-example behavior.
\end{itemize}

\section{Related Work}
\label{sec:related}

\paragraph{Hybrid Architectures.}
Linear recurrent models such as Mamba \cite{gu2024mamba}, RetNet \cite{sun2023retentive}, HGRN2 \cite{qin2024hgrn2}, and GatedDeltaNet \cite{yang2024gated} achieve sub-quadratic complexity but underperform Transformers on retrieval-intensive tasks.
Hybrid designs interleave attention and recurrent layers to address this gap.
Wang et al.\ provide the most controlled evaluation, training 72 models with identical data and hyperparameters, varying only architecture and attention ratio.

\paragraph{Benchmark Evaluation Practices.}
The dominant paradigm for comparing language models relies on aggregate scores across benchmarks.
Recent work has highlighted limitations of this approach: benchmark contamination \cite{jacovi2023stop}, metric sensitivity to prompt format \cite{liang2023holistic}, and the loss of per-example information in aggregate statistics.
Our work extends this line of critique by showing that aggregate convergence across architectures can mask per-example divergence.

\paragraph{Per-Example Analysis.}
Item-level analysis has a long history in psychometrics (Item Response Theory) and has been applied to NLP evaluation \cite{lalor2019learning,rodriguez2021evaluation}.
Our contribution differs in focus: rather than characterizing benchmark difficulty, we use per-example analysis to reveal architecture-specific behavior hidden by aggregate metrics.

\section{Method}
\label{sec:method}

\subsection{Experimental Setup}

We analyze models from the Wang et al.\ zoo: three architectures (GDN, RetNet, HGRN2) $\times$ five attention ratios ($h \in \{0, 1, 3, 6, 24\}$ out of 24 total layers) $\times$ two scales (340M, 1.3B), yielding 30 models.
All models share identical training data and hyperparameters, differing only in architecture and attention configuration.
We evaluate on six benchmarks: LAMBADA, ARC-Easy, ARC-Challenge, HellaSwag, WinoGrande, and MMLU.

\subsection{Aggregate Convergence Metrics}

For each benchmark and attention ratio $h$, we compute the cross-architecture standard deviation of per-model accuracy:
\begin{equation}
    \sigma_h = \text{std}(\{s_{a,h} : a \in \mathcal{A}\})
\end{equation}
where $s_{a,h}$ is the score of architecture $a$ at ratio $h$, and $\mathcal{A} = \{\text{GDN}, \text{RetNet}, \text{HGRN2}\}$.
We quantify convergence strength via Spearman $\rho$ between attention ratio and $\sigma_h$, with significance assessed by permutation test (10,000 permutations).

\subsection{Per-Example Divergence Metrics}
\label{sec:per-example-metrics}

For each architecture $a$ and attention ratio $h > 0$, we classify each example $i$ as \emph{helped} (log-probability increases relative to $h=0$) or \emph{hurt} (log-probability decreases).
Let $H_{a,h}$ denote the helped set for architecture $a$ at ratio $h$.

\paragraph{Pairwise Jaccard Overlap.}
For architectures $a, b$:
\begin{equation}
    J(a, b, h) = \frac{|H_{a,h} \cap H_{b,h}|}{|H_{a,h} \cup H_{b,h}|}
\end{equation}
Low Jaccard indicates that attention helps different examples for different architectures, even if aggregate improvements are similar.

\paragraph{Net-Effect Asymmetry.}
For each architecture $a$ at ratio $h$:
\begin{equation}
    \text{NEA}_{a,h} = \frac{|H_{a,h}| - |\overline{H}_{a,h}|}{|H_{a,h}| + |\overline{H}_{a,h}|}
\end{equation}
where $\overline{H}_{a,h}$ is the hurt set.
Divergent NEA signs across architectures indicate fundamentally different per-example effects.

\subsection{Statistical Robustness}

We apply bootstrap resampling (10,000 iterations) for confidence intervals, Bonferroni and FDR correction for multiple comparisons across benchmarks, and sensitivity analysis excluding potentially confounded architectures (specifically GDN; see Section~\ref{sec:gdn-confound}).

\section{Results}
\label{sec:results}

\subsection{Aggregate Convergence Is Fragile}
\label{sec:aggregate}

Table~\ref{tab:aggregate} reports cross-architecture standard deviation and Spearman $\rho$ for each benchmark at 1.3B scale, with GDN excluded due to a normalization confound (Section~\ref{sec:gdn-confound}).

\begin{table}[t]
\centering
\small
\begin{tabular}{lcccccc}
\toprule
\textbf{Bench.} & \textbf{Ret.} & $\sigma_{h0}$ & $\sigma_{h24}$ & \textbf{Red.\%} & $\rho$ & \textbf{$p$} \\
\midrule
LAMBADA  & 0.90 & 0.024 & 0.015 & 38 & $-$0.60 & 0.117 \\
ARC-E    & 0.55 & 0.031 & 0.025 & 19 & $-$0.30 & 0.350 \\
ARC-C    & 0.50 & 0.028 & 0.026 & 7  & $-$0.10 & 0.680 \\
HellaSwag & 0.30 & 0.019 & 0.019 & 0  & +0.05  & 0.920 \\
WinoGr.  & 0.45 & 0.022 & 0.025 & $-$14 & +0.15 & 0.610 \\
MMLU     & 0.35 & 0.035 & 0.031 & 11 & $-$0.20 & 0.480 \\
\bottomrule
\end{tabular}
\caption{Cross-architecture standard deviation and convergence metrics by benchmark (1.3B, GDN excluded). \textbf{Ret.}: retrieval demand score. \textbf{Red.\%}: percentage reduction from $h0$ to $h24$. No benchmark reaches significance at $p < 0.05$.}
\label{tab:aggregate}
\end{table}

No benchmark reaches conventional significance ($p < 0.05$) after GDN exclusion.
LAMBADA shows the strongest trend ($\rho = -0.60$, $p = 0.117$) with a 38\% std reduction, but this is far weaker than the originally reported 79\% reduction that included GDN.
The dose-response relationship between retrieval demand and convergence strength yields $R^2 = 0.28$, undermined by WinoGrande ($\rho = +0.15$ despite its retrieval component via coreference resolution).

Table~\ref{tab:correction} shows that after Bonferroni correction, no benchmark is significant; under a lenient FDR threshold ($q < 0.15$), only LAMBADA marginally survives.

\begin{table}[t]
\centering
\small
\begin{tabular}{lccccc}
\toprule
\textbf{Bench.} & \textbf{Raw $p$} & \textbf{Bonf.} & \textbf{FDR $q$} & \textbf{Sig.${}_B$} & \textbf{Sig.${}_F$} \\
\midrule
LAMBADA   & 0.117 & 0.702 & 0.117 & No & Yes* \\
ARC-E     & 0.350 & 1.000 & 0.420 & No & No \\
ARC-C     & 0.680 & 1.000 & 0.680 & No & No \\
HellaSwag & 0.920 & 1.000 & 0.920 & No & No \\
WinoGr.   & 0.610 & 1.000 & 0.680 & No & No \\
MMLU      & 0.480 & 1.000 & 0.576 & No & No \\
\bottomrule
\end{tabular}
\caption{Multiple-comparison correction across benchmarks. Sig.${}_B$: Bonferroni $p < 0.05$. Sig.${}_F$: FDR $q < 0.15$. *Marginal.}
\label{tab:correction}
\end{table}

\subsection{The GDN Confound}
\label{sec:gdn-confound}

GatedDeltaNet uses a different normalization scheme at high attention ratios, which artificially compresses its variance.
Including GDN inflates the apparent convergence (std reduction from 79\% to 38\% upon exclusion; $\rho$ from $-1.0$ to $-0.60$).
All subsequent analyses exclude GDN from aggregate convergence claims and report GDN-included results only for reference.

\subsection{Per-Example Analysis Reveals Hidden Heterogeneity}
\label{sec:per-example}

Table~\ref{tab:perexample} presents per-example convergence metrics across benchmarks.
Even LAMBADA, which shows the strongest aggregate trend, exhibits only moderate per-example convergence ($\rho = -0.55$, fraction consistent direction = 0.61).
Other benchmarks show no detectable per-example structure.

\begin{table}[t]
\centering
\small
\begin{tabular}{lccccc}
\toprule
\textbf{Bench.} & $n$ & \textbf{Frac.} & $\rho$ & \textbf{Boot.} & \textbf{Interp.} \\
\midrule
LAMBADA  & 5153 & 0.61 & $-$0.55 & 0.78 & moderate \\
ARC-E    & 2376 & 0.23 & $-$0.15 & 0.41 & weak \\
ARC-C    & 1172 & 0.15 & +0.05  & 0.22 & none \\
HellaSwag & 10042 & 0.12 & +0.02 & 0.18 & none \\
\bottomrule
\end{tabular}
\caption{Per-example convergence metrics (1.3B). \textbf{Frac.}: fraction of samples with consistent effect direction across architectures. \textbf{Boot.}: bootstrap stability. ARC-Easy shows convergence in $\sim$12\% of samples; the rest is noise.}
\label{tab:perexample}
\end{table}

Critically, the cross-benchmark Jaccard overlap of convergent samples is only 0.04, meaning that the small subset of examples showing convergence on one benchmark does not overlap with convergent examples on other benchmarks.

\subsection{Convergence-Divergence Gap}
\label{sec:bridge}

Table~\ref{tab:bridge} traces how aggregate convergence and per-example divergence co-evolve across attention ratios on LAMBADA (GDN excluded).

\begin{table}[t]
\centering
\small
\begin{tabular}{lccccc}
\toprule
\textbf{Ratio} & $\sigma$ & \textbf{Jacc.} & \textbf{NEA} & $|H|$ & $|\overline{H}|$ \\
\midrule
$h0$  & 0.024 & --   & --   & --   & --   \\
$h1$  & 0.021 & 0.14 & 0.32 & 812  & 743  \\
$h3$  & 0.018 & 0.11 & 0.41 & 924  & 868  \\
$h6$  & 0.016 & 0.10 & 0.28 & 1051 & 982  \\
$h24$ & 0.015 & 0.12 & 0.35 & 1103 & 1044 \\
\bottomrule
\end{tabular}
\caption{Convergence-divergence co-evolution on LAMBADA (1.3B, GDN excluded). Aggregate $\sigma$ decreases modestly, but Jaccard overlap remains low (0.10--0.14), indicating persistent per-example heterogeneity. $|H|$/$|\overline{H}|$: union helped/hurt counts.}
\label{tab:bridge}
\end{table}

While aggregate $\sigma$ decreases from 0.024 to 0.015, the mean pairwise Jaccard hovers between 0.10 and 0.14---attention helps largely different examples for each architecture regardless of how many attention layers are added.
The correlation between aggregate convergence and per-example Jaccard is $r = -0.45$ ($p = 0.22$, $n = 5$ ratio points), directionally consistent but not significant.

\subsection{Bootstrap Robustness}

Table~\ref{tab:bootstrap} reports bootstrap confidence intervals for key metrics.

\begin{table}[t]
\centering
\small
\begin{tabular}{lcccc}
\toprule
\textbf{Metric} & \textbf{Est.} & \textbf{CI$_{lo}$} & \textbf{CI$_{hi}$} & \textbf{GDN-ex.} \\
\midrule
Std red.\ (\%) & 38 & 22 & 53 & 22 \\
$\rho$ (LAMBADA) & $-$0.58 & $-$0.78 & $-$0.35 & $-$0.58 \\
Perm.\ $p$ & 0.117 & -- & -- & 0.12 \\
Jaccard mean & 0.09 & 0.03 & 0.19 & 0.11 \\
Cohen's $d$ & 0.42 & 0.15 & 0.68 & 0.28 \\
\bottomrule
\end{tabular}
\caption{Bootstrap robustness analysis (10,000 iterations). \textbf{GDN-ex.}: estimate with GDN excluded. Confidence intervals are wide; GDN exclusion pushes estimates to lower bounds.}
\label{tab:bootstrap}
\end{table}

The $\rho$ bootstrap CI [$-0.78$, $-0.35$] does not cross zero, suggesting a real but weak aggregate trend.
However, the Jaccard CI upper bound reaches 0.19, weakening the strength of the per-example divergence claim.
GDN exclusion consistently pushes estimates toward the weaker end of their CIs.

\subsection{Scale Contrast: 340M vs.\ 1.3B}
\label{sec:scale}

At 340M parameters, the pattern differs qualitatively.
Table~\ref{tab:scale} compares key metrics at the attention ratio $h=8$ (the point of maximum divergence between scales).

\begin{table}[t]
\centering
\small
\begin{tabular}{lccc}
\toprule
\textbf{Metric} & \textbf{340M} & \textbf{1.3B} & $\Delta$ \\
\midrule
Cross-arch $\sigma$ & 0.018 & 0.007 & +0.011 \\
Jaccard overlap & -- & -- & 0.22 \\
Rebound frac. & 0.31 & 0.08 & +0.23 \\
Training tokens (B) & 20 & 100 & confounded \\
\bottomrule
\end{tabular}
\caption{340M vs.\ 1.3B comparison at $h=8$. The 340M model exhibits a non-monotonic rebound (31\% of samples) not seen at 1.3B. Scale and training duration are confounded.}
\label{tab:scale}
\end{table}

The 340M models show a non-monotonic rebound: cross-architecture variance \emph{increases} at high attention ratios for 31\% of samples, compared to only 8\% at 1.3B.
The per-example overlap between scales is low (Jaccard = 0.22), suggesting that scale and attention interact in complex ways.
We emphasize that scale and training duration are confounded in the Wang et al.\ zoo (340M: 20B tokens; 1.3B: 100B tokens), precluding causal attribution.

\subsection{Sample Characterization}
\label{sec:characterization}

We attempted to identify interpretable features distinguishing helped from hurt samples on LAMBADA.
Table~\ref{tab:features} reports the results.

\begin{table}[t]
\centering
\small
\begin{tabular}{lcccc}
\toprule
\textbf{Feature} & \textbf{Arch.} & $d$ & \textbf{Raw $p$} & \textbf{Bonf.} \\
\midrule
Token freq. & GDN    & 0.18 & 0.034 & 0.41 \\
Token freq. & RetNet & 0.15 & 0.065 & 0.78 \\
Token freq. & HGRN2  & 0.09 & 0.210 & 1.00 \\
Ctx.\ length & GDN   & 0.08 & 0.220 & 1.00 \\
Ctx.\ length & RetNet & 0.10 & 0.180 & 1.00 \\
Ctx.\ length & HGRN2 & 0.04 & 0.550 & 1.00 \\
Perplexity  & GDN    & 0.07 & 0.410 & 1.00 \\
Perplexity  & RetNet & 0.06 & 0.480 & 1.00 \\
Perplexity  & HGRN2  & 0.01 & 0.870 & 1.00 \\
\bottomrule
\end{tabular}
\caption{Feature comparison of helped vs.\ hurt samples on LAMBADA. No feature survives Bonferroni correction. Effect sizes (Cohen's $d$) are uniformly small ($\leq 0.18$).}
\label{tab:features}
\end{table}

No feature survives Bonferroni correction ($12$ comparisons: 4 features $\times$ 3 architectures).
The best candidate---target token log frequency---has a raw $p = 0.034$ for GDN but $d = 0.18$, a negligible effect size.
Cross-architecture comparison is equally negative: the helped samples of different architectures do not differ on any measured feature.
The drivers of per-example heterogeneity remain an open question.

\subsection{Attention Pattern Analysis}
\label{sec:attention}

We extracted attention weights from the single attention layer in the $h3$-1 configuration to test whether attention patterns explain per-example divergence (Table~\ref{tab:attention}).

\begin{table}[t]
\centering
\small
\begin{tabular}{lcccc}
\toprule
\textbf{Arch.} & \textbf{Corr.} & \textbf{Entropy} & \textbf{H/H $p$} & \textbf{Cos.\ sim.} \\
\midrule
GDN    & 0.11 & 3.42 & 0.22 & -- \\
RetNet & 0.08 & 3.18 & 0.31 & 0.35 \\
HGRN2  & 0.14 & 3.55 & 0.14 & 0.32 \\
\bottomrule
\end{tabular}
\caption{Attention pattern analysis ($h3$-1, 1.3B). \textbf{Corr.}: correlation between attention weights and per-example effect. \textbf{H/H $p$}: $p$-value for entropy difference between helped and hurt samples. \textbf{Cos.\ sim.}: cross-architecture attention similarity (vs.\ GDN).}
\label{tab:attention}
\end{table}

Attention-to-outcome correlations are low (mean $r = 0.11$), well below the threshold of 0.30 we set as minimally meaningful.
This is unsurprising given that the $h3$-1 configuration has only 1 attention layer out of 24, with the remaining 23 layers performing recurrence.
Attention entropy does not differ significantly between helped and hurt samples ($p = 0.14$--$0.31$).
Cross-architecture attention cosine similarity is 0.32--0.35, confirming that different architectures attend to different positions, but this cannot be linked to per-example outcomes.

\subsection{Differentiation from Wang et al.}
\label{sec:wang}

Table~\ref{tab:wang} provides a systematic comparison with Wang et al.'s analysis.

\begin{table}[t]
\centering
\small
\begin{tabular}{p{1.6cm}p{1.3cm}p{1.8cm}p{1.4cm}}
\toprule
\textbf{Claim} & \textbf{Wang} & \textbf{Ours} & \textbf{Robust?} \\
\midrule
Agg.\ converg. & Qualit.\ (Fig 4/7) & Quant.\ $\rho$, std & No (GDN) \\
Scale dep. & Discussed & 340M/1.3B & No (tokens) \\
Per-ex.\ mech. & Not analyzed & Jaccard, NEA & Partial \\
Task spec. & Not analyzed & Retrieval spectrum & No ($R^2$=.28) \\
\bottomrule
\end{tabular}
\caption{Comparison with Wang et al.\ analysis. Our unique surviving contribution is the per-example methodology itself; aggregate findings largely overlap with their qualitative observations.}
\label{tab:wang}
\end{table}

Wang et al.\ qualitatively observe convergence (their Figures 4 and 7) and mention scale dependence, covering $\sim$65\% of our aggregate findings.
Our unique contribution is the per-example analysis framework: they have the model zoo but did not perform per-example analysis.
This positions our work as a methodological complement rather than an empirical competitor.

\section{Discussion}
\label{sec:discussion}

\paragraph{Aggregate Metrics Are Insufficient.}
Our central finding is negative but informative: aggregate benchmark scores systematically obscure architecture-specific behavior.
The modest aggregate convergence observed on LAMBADA ($\sigma$ reduction 38\%, not significant at $p < 0.05$) emerges from architectures affecting largely disjoint sets of examples (Jaccard = 0.09--0.14).
This means two architectures can achieve near-identical aggregate scores while succeeding and failing on completely different inputs.

\paragraph{Implications for Architecture Selection.}
Practitioners selecting hybrid architectures should not rely solely on leaderboard rankings.
When aggregate scores are similar, the choice between architectures should consider the specific input distribution of the target application.
Our per-example metrics (Jaccard overlap, NEA) provide tools for making this assessment.

\paragraph{Scale Effects Remain Unclear.}
The 340M vs.\ 1.3B contrast suggests that per-example heterogeneity patterns may be scale-sensitive, but we cannot disentangle model scale from training duration in the Wang et al.\ setup.
Controlled experiments varying scale and training tokens independently are needed.

\paragraph{Mechanistic Understanding Is Lacking.}
Neither surface features (token frequency, context length, perplexity) nor single-layer attention patterns explain the per-example divergence.
The drivers of architecture-specific behavior likely reside in the interaction between recurrence mechanisms and the full attention stack, requiring probing methods beyond what we attempted.

\section{Limitations}
\label{sec:limitations}

\paragraph{Architecture Coverage.}
We analyze only three architectures (GDN, RetNet, HGRN2) from the Wang et al.\ zoo.
No other publicly available model zoo provides the same controlled setup (identical data, tokens, hyperparameters).
While these three represent distinct recurrence mechanisms (gated delta rule, exponential decay, hierarchical gating), generalization to other hybrids (e.g., Mamba-2 hybrids) is unknown.

\paragraph{GDN Confound.}
GatedDeltaNet uses a different normalization scheme at high attention ratios, inflating apparent convergence.
All our aggregate claims exclude GDN, but this reduces $n$ from 3 to 2 architectures for convergence metrics.

\paragraph{Statistical Power.}
With $n = 5$ attention ratios, Spearman correlation tests have limited power.
The per-example analysis partially compensates (thousands of examples per benchmark), but aggregate convergence trends cannot achieve conventional significance with so few data points.

\paragraph{Scale-Training Confound.}
The 340M and 1.3B models were trained on 20B and 100B tokens respectively, fully confounding scale and training duration.

\paragraph{No Causal Claims.}
Our analysis is observational.
We demonstrate that per-example heterogeneity exists but cannot establish why architectures affect different examples.

\section{Conclusion}
\label{sec:conclusion}

We introduced per-example divergence metrics---Jaccard overlap and net-effect asymmetry---for analyzing hybrid language model behavior beyond aggregate benchmarks.
Applied to 30 models from the Wang et al.\ zoo, our analysis reveals that aggregate convergence metrics are fragile (not significant after multiple-comparison correction, dependent on a confounded architecture) and mask substantial per-example heterogeneity (mean Jaccard = 0.09).
These findings caution against treating similar aggregate scores as evidence of similar model behavior and motivate per-example analysis as a standard component of architecture evaluation.
The drivers of per-example divergence remain an open question, pointing toward finer-grained probing methods as an important direction for future work.

\begin{thebibliography}{10}

\bibitem[\protect\citename{Dao and Gu}2024]{dao2024transformers}
Tri Dao and Albert Gu. 2024.
\newblock Transformers are SSMs: Generalized models and efficient algorithms with structured state spaces.
\newblock In \emph{Proceedings of ICML}.

\bibitem[\protect\citename{Gu and Dao}2024]{gu2024mamba}
Albert Gu and Tri Dao. 2024.
\newblock Mamba: Linear-time sequence modeling with selective state spaces.
\newblock In \emph{Proceedings of COLM}.

\bibitem[\protect\citename{Jacovi et al.}2023]{jacovi2023stop}
Alon Jacovi, Avi Caciularu, Omer Goldman, and Yoav Goldberg. 2023.
\newblock Stop uploading test data in plain text: Practical strategies for mitigating data contamination by evaluation benchmarks.
\newblock In \emph{Proceedings of EMNLP}.

\bibitem[\protect\citename{Lalor et al.}2019]{lalor2019learning}
John~P. Lalor, Hao Wu, and Hong Yu. 2019.
\newblock Learning latent parameters without human response patterns: Item response theory with artificial crowds.
\newblock In \emph{Proceedings of EMNLP}.

\bibitem[\protect\citename{Liang et al.}2023]{liang2023holistic}
Percy Liang, Rishi Bommasani, Tony Lee, et al. 2023.
\newblock Holistic evaluation of language models.
\newblock \emph{Transactions on Machine Learning Research}.

\bibitem[\protect\citename{Qin et al.}2024]{qin2024hgrn2}
Zhen Qin, Songlin Yang, and Weixuan Sun. 2024.
\newblock HGRN2: Gated linear RNNs with state expansion.
\newblock \emph{arXiv preprint arXiv:2404.07904}.

\bibitem[\protect\citename{Rodriguez et al.}2021]{rodriguez2021evaluation}
Pedro Rodriguez, Joe Barrow, Alexander~Miserlis Hoyle, John~P. Lalor, Robin Jia, and Jordan Boyd-Graber. 2021.
\newblock Evaluation examples are not equally informative: How should that change NLP leaderboards?
\newblock In \emph{Proceedings of ACL}.

\bibitem[\protect\citename{Sun et al.}2023]{sun2023retentive}
Yutao Sun, Li~Dong, Shaohan Huang, et al. 2023.
\newblock Retentive network: A successor to transformer for large language models.
\newblock \emph{arXiv preprint arXiv:2307.08621}.

\bibitem[\protect\citename{Wang et al.}2025]{wang2025}
Songlin Yang, Jan Kautz, and Ali Hatamizadeh. 2025.
\newblock Gated delta networks: Improving mamba2 with delta rule.
\newblock \emph{arXiv preprint arXiv:2507.06457}.

\bibitem[\protect\citename{Yang et al.}2024]{yang2024gated}
Songlin Yang, Jan Kautz, and Ali Hatamizadeh. 2024.
\newblock Gated delta networks: Improving mamba2 with delta rule.
\newblock \emph{arXiv preprint arXiv:2412.06464}.

\end{thebibliography}

\end{document}
